% 字を小さくした箇条書き・用語定義・引用。**本文の誤削除**を測る標本。
%
% 逐語ブロック（lstlisting / verbatim）の判定を「本文と違う字の大きさ × 文章として
% 読めない × 4 行以上 × どの行も文末記号で終わらない × 右端がばらける」で決めていた頃、
% **句点を打たない短い箇条書きがこれを全部満たして**丸ごと「本文以外（訳しません）」に
% 落ちていた。貢献一覧・用語定義・引用は、読者が訳をいちばん欲しがる本文。
%
% 直す前の実測: `{\small itemize}` の 5 項目、`\footnotesize description` の用語定義
% 5 件、`{\small quote}` の 5 行が、それぞれまとめて付物になった。しかも逐語ブロックは
% 定義上「文章として読めない行」だけで出来ているので、**採録率は原理的に 1 のまま**で
% 申告も出ない（`coverage: 1` / `readingOrderRisk: low` / 継ぎ目の印なし）。
%
% 見分けるのは字の大きさではなく**記号の有無**。コードは括弧・等号・コロン・下線が
% 地の文になるが、箇条書きも用語定義も引用もそれを 1 つも持たない（layout.ts の
% CODE_SYMBOLS を参照）。
%
% 句読点が無いので、塊の中の行は 1 文につながる。そこが正解（`\sn{...}`）。
% `\sn{...}` は「本文 1 文だが、版面はすぐ下に別に書く」印で、組版はしない。
% 中身は、下の塊の行を版面の順に空白でつないだもの。
% 塊の直後に節見出しを置いてあるのは、塊が次の段落と 1 文になると
% 「本文として残ったか」だけを測れなくなるため。
\documentclass[10pt]{article}

\usepackage{mathptmx}
\usepackage[margin=2cm]{geometry}

\newcommand{\s}[1]{#1}
\newcommand{\sn}[1]{}

\title{Notes on Block Selection}
\author{Anonymous}
\date{}

\begin{document}
\maketitle

\section{Introduction}

\s{Sparse attention reduces the quadratic cost of long context models to a manageable level.}
\s{Our contributions are summarised in the following list.}

\sn{• A block sparse attention kernel • An analysis of the memory traffic • A scheduler for mixed block sizes • Experiments on six public corpora • An open source reference implementation}
{\small
\begin{itemize}
\item A block sparse attention kernel
\item An analysis of the memory traffic
\item A scheduler for mixed block sizes
\item Experiments on six public corpora
\item An open source reference implementation
\end{itemize}
}

\section{Notation}

\s{We use the following notation throughout the paper.}

\sn{Block a contiguous span of keys Pattern the set of blocks that a query attends to Budget the number of blocks per query Router the module that selects blocks Kernel the fused implementation on the accelerator}
{\footnotesize
\begin{description}
\item[Block] a contiguous span of keys
\item[Pattern] the set of blocks that a query attends to
\item[Budget] the number of blocks per query
\item[Router] the module that selects blocks
\item[Kernel] the fused implementation on the accelerator
\end{description}
}

\section{Related work}

\s{The original proposal is stated in the following terms, which we quote in full.}

\sn{attention is all you need and memory is all you have so spend it where the mass is and nowhere else at all that is the whole of the method}
{\small
\begin{quote}
attention is all you need\\
and memory is all you have\\
so spend it where the mass is\\
and nowhere else at all\\
that is the whole of the method
\end{quote}
}

\section{Results}

\s{The sparse kernel reaches the same accuracy as the dense baseline on every task that we tried.}

\end{document}
